\section*{ID9311: Definition of T\_mathfrakheld}
\paragraph{Metadata.}
PiID: 8589638140918 | RTEID: RTE:P36156.6808:T3.0264 | UUID: 62373df8-9cab-52b1-b309-904bc22fb201

\paragraph{Canonical Statement.}
\$T\_\{\textbackslash{}mathfrak\{h\}eld\}=\textbackslash{}frac\{\textbackslash{}epsilon\_\{0\}\}\{2\}\textbackslash{}int\_\{\textbackslash{}mathbb\{R\}\textasciicircum{}\{3\}\}(\textbackslash{}frac\{\textbackslash{}partial\textbackslash{}vec\{A\}\}\{\textbackslash{}partial t\})\textbackslash{}cdot(\textbackslash{}frac\{\textbackslash{}partial\textbackslash{}vec\{A\}\}\{\textbackslash{}partial t\})d\textasciicircum{}\{3\}x\$

\paragraph{Canonical Equation.}
\[
T_{\mathfrak{h}eld}=\frac{\epsilon_{0}}{2}\int_{\mathbb{R}^{3}}(\frac{\partial\vec{A}}{\partial t})\cdot(\frac{\partial\vec{A}}{\partial t})d^{3}x
\]

\paragraph{Dictionary.}
Definition of T\_mathfrakheld — T\_mathfrakheld=fracε\_02∫\_R\textasciicircum{}3((partialA)/(∂ t))·((partialA)/(∂ t))d\textasciicircum{}3x

\paragraph{Encyclopedia.}
Definition of T\_mathfrakheld is a variational / legendre derivative, written T\_mathfrakheld=fracε\_02∫\_R\textasciicircum{}3((partialA)/(∂ t))·((partialA)/(∂ t))d\textasciicircum{}3x. It arises from a variational or Legendre procedure, defining conjugate quantities or stationarity conditions. It is built from ∂, A, x, defining T\_mathfrakheld. Within the Trawin Zero Point registry it serves as one element of the framework's mathematical narrative.
